% !TEX program = xelatex
\documentclass[a4paper,12pt]{article}

\usepackage{fontspec}
\usepackage{xeCJK}
\setmainfont{Times New Roman}
\setCJKmainfont{IPAexMincho}

\usepackage[margin=25mm]{geometry}
\usepackage{enumitem}

\title{生成AI利用申告書}
\date{\today}

\begin{document}

\maketitle

\vspace{1em}

\noindent
\begin{itemize}
  \item 氏名：
  \item 学籍番号：
  \item プログラムファイル名：
  \item 使用したAIツール（例：ChatGPT 等）：
\end{itemize}

\vspace{1em}

\section*{1. 使用目的}
%本課題において、生成AIをどのような用途で利用したかについて、具体的に記述してください。\\
%（例：エラーの修正、コードの作成、アイデアの検討 等）

\section*{2. AIへの指示（プロンプト）}
%生成AIに対して与えた指示または質問の内容について、主要なものを1〜2点記述してください。

\section*{3. AI出力の利用方法}
\begin{itemize}
  \item □ そのまま利用した
  \item □ 一部修正のうえ利用した
  \item □ 参考情報としてのみ利用した
\end{itemize}
%上記の選択について、その理由を具体的に記述してください。

\section*{4. 正当性の確認方法}
%生成AIの出力内容が正しいものであるかについて、どのような方法により確認を行ったかを記述してください。
%（例）
%\begin{itemize}
%  \item 実行結果により動作確認を行った
%  \item 異なる入力条件により検証を行った
%  \item 内容を精読し、意味を理解した
%\end{itemize}

%確認方法については、具体的に記述すること。

\section*{5. 問題点および誤りの有無}

\begin{itemize}
  \item □ 問題または誤りがあった
  \item □ 問題または誤りは確認されなかった
\end{itemize}
%上記の確認結果に基づき、検証内容および気付いた点について具体的に記述してください。

\section*{6. 独自に工夫した点}
%最終的な成果物において、自身で考え工夫した点について具体的に記述してください。

\section*{参考文献}

\end{document}